% =========================================================================
% Writing the robotics paper -- a craft guide.
%
% Section-by-section guidance on what each part of a paper must do and how
% to write it: the job each section performs, what a good sentence in it
% looks like, and the habits that separate a paper reviewers argue for from
% one they merely fail to reject.
%
% Compile:  pdflatex award_paper_guide.tex   (twice, for the ToC)
% =========================================================================
\documentclass[11pt,a4paper]{article}

\usepackage[margin=2.4cm]{geometry}
% lmodern BEFORE fontenc: T1 with plain Computer Modern sends MiKTeX to
% METAFONT to render bitmap EC fonts, which takes minutes. Latin Modern
% ships scalable Type-1 faces for T1 and compiles instantly.
\usepackage{lmodern}
\usepackage[T1]{fontenc}
\usepackage{amssymb}
\usepackage{booktabs}
\usepackage{tabularx}
\usepackage{enumitem}
\usepackage{xcolor}
\usepackage[hidelinks]{hyperref}

\definecolor{crimson}{HTML}{9D2235}
\hypersetup{colorlinks=true,linkcolor=crimson,urlcolor=crimson}

\newcommand{\lead}[1]{\medskip\noindent\textbf{#1}\quad}
\newcommand{\ck}{\item[$\square$]}
\setlist[itemize]{itemsep=3pt,parsep=0pt,topsep=4pt}
\setlist[enumerate]{itemsep=3pt,parsep=0pt,topsep=4pt}

\title{\vspace{-1.4cm}Writing the Robotics Paper\\
  \large What each section must do, and how to write it}
\author{}
\date{}

\begin{document}
\maketitle
\vspace{-1.8cm}

\begin{abstract}
\noindent Section-by-section craft guidance: what job each part of a paper
performs, what a good sentence in it looks like, and the specific habits that
separate a paper reviewers argue for from one they merely fail to reject.
Written to be used on a draft that already exists.
\end{abstract}

\tableofcontents
\vspace{0.4cm}

% =========================================================================
\section{The governing idea}

A paper defends \emph{one} claim. Everything else --- the method, the three
experiments, the ablations --- is a witness to that claim. Before revising
anything, write the claim as a single sentence that could be false. If you
cannot, the paper has a structural problem no amount of editing will fix.

The test for every paragraph thereafter is: \emph{which part of that sentence
does this support?} Material that supports no part of it is why papers run
over length.

\lead{Coin the claim, then never paraphrase it.} Pick the words once and use
exactly those words in the title if possible, the abstract, the section that
establishes the claim, and the conclusion. Drafts routinely carry three
near-synonymous phrasings of their own thesis, which gives the reader three
chances to remember none of them. A reader who can repeat your sentence at a
poster session is a reader who will cite you.

% =========================================================================
\section{Before you draft}

The most common failure in a first paper is not a weak abstract; it is
writing up work that does not yet contain a claim. Apply the test above
\emph{before} opening the editor: can you complete the sentence ``we show
that $X$, and the evidence is $Y$'' with $Y$ already on disk? If the
evidence is not collected, you are still designing experiments, and writing
prose now will only disguise that.

\lead{Choose the claim your strongest evidence supports.} A body of results
usually admits several candidate claims. Pick the one your best evidence
proves, not the most ambitious one you could almost defend --- the ambitious
version becomes the future-work paragraph or the next paper. The choice is
consequential: the claim determines which baselines you owe, which
experiments are load-bearing, and which results are merely interesting.

\lead{Build the figures and tables first.} Before writing prose, assemble
the tables and figures as they will appear --- real data, final columns,
captions drafted --- and pin them to a skeleton of section headings. If the
figures do not carry the argument on their own, more writing will not save
them, and it is far cheaper to discover a missing experiment at this stage
than after the text is set around its absence. Then draft the abstract as a
plan of what the paper will claim, knowing you will rewrite it last.

\lead{Mixed results can still be a paper.} If the method wins somewhere and
loses somewhere, the paper about the \emph{boundary} --- where it holds,
where it breaks, and why --- is often stronger than a strained paper about
the win. Characterisation is a contribution; torturing the framing until the
loss disappears is how papers get rejected by the one reviewer who reruns
the numbers.

% =========================================================================
\section{Title}

Eight words is the working target. Two constructions carry most good titles:
a short pronounceable system name, a colon, then the mechanism; or a claim
stated as a sentence with an active verb and no colon at all.

\begin{itemize}
  \item \textbf{One claim, not three.} Titles that enumerate contributions
    read as a table of contents. Pick the one that would make a reader stop.
  \item \textbf{Name the mechanism, not the field.} A title that could sit on
    twenty other papers is a wasted asset.
  \item \textbf{Banned openers:} ``Towards'', ``A Novel'', ``A Framework
    for'', ``On the''. They signal hedging before the reader has read a word.
  \item \textbf{No numerals}, and no acronym you have not earned the right to
    coin. If you do coin one, check it does not already mean something else in
    the field.
\end{itemize}

% =========================================================================
\section{Abstract}

\textbf{Every sentence carries weight.} The abstract is the only part of the
paper most people read, and it is short enough that a reader notices a wasted
line. The test: delete any sentence and ask whether information is lost. If
not, it was scaffolding --- ``In this paper, we\ldots'', ``Extensive
experiments were conducted\ldots'' --- and it should go.

Seven sentences, one job each:

\begin{enumerate}[leftmargin=*,label=\textbf{S\arabic*.}]
  \item \textbf{The enemy.} A concrete failure of how things are currently
    done, stated as a property rather than a statistic. Not ``robotics is
    important'', but the specific thing that breaks.
  \item \textbf{Why the obvious fix fails.} Credit the alternative a reader is
    already thinking of, then name the mechanism by which it fails. This
    sentence is where a citation earns its place.
  \item \textbf{What you did.} ``We present X'', with the mechanism inside
    the same sentence, phrased as a contrast: \emph{instead of} A, we do B.
    A reader should be able to restate your method from this sentence alone.
  \item \textbf{The central claim, with its number.} One number. The one that
    would survive if the paper were compressed to a tweet.
  \item \textbf{What you tested it on, and how you made it hard.} Name the
    datasets and the baselines --- and say how you strengthened the
    opposition: reproduced their published result first, matched operating
    points, handed them privileged information.
  \item \textbf{The result, with its boundary condition.} Give the win and
    the axis it was traded on --- the scope within which the claim holds ---
    rather than the bare deficit. ``Within $X$ of the dense baseline at a
    thousandth the memory, and ahead of it in every degraded condition
    tested'' is honest, quotable, and standard. A \emph{lost comparison}
    belongs in the results table with a diagnosis, not in the abstract:
    reviewers have little precedent for reading an abstract-level concession
    charitably, and it displaces the sentence that should be selling the
    work.
  \item \textbf{The thesis, coined.} The sentence you want quoted.
\end{enumerate}

\lead{One number, not five.} Scattering metrics through an abstract is the
single most common tell of an ordinary paper. Keep the headline figure and,
if the object has an identity number (a size, a rate, a count), keep that.
Everything else moves to the first figure's caption or the results tables.

\lead{Counts are underused.} ``In all nine conditions tested'' asserts
coverage in a way a mean cannot, and is much harder to argue with than an
average. If you have run the sweep, say how many.

\lead{Do not end on effectiveness.} ``Experiments demonstrate the
effectiveness of the proposed method'' is content-free and it is where a
reader's attention goes to die. End on the claim.

% =========================================================================
\section{Introduction}

The introduction expands the abstract's seven moves into three to six
paragraphs and ends with the contributions. Its job is to make the reader
believe the problem is real and that you have understood it precisely.

\begin{itemize}
  \item \textbf{Reach the specific gap within three or four sentences.} A
    general opening about the importance of the field costs you the reader you
    most want.
  \item \textbf{Do not open on a statistic.} Open on the failure.
  \item \textbf{Credit before you displace.} Describe the standard approach
    as a competent person would describe it, then show precisely where it
    breaks. A strawman tells the reviewer --- who may have written the
    strawman --- that you did not read carefully.
  \item \textbf{The gap must be falsifiable.} ``Little work has been done''
    is not a gap; it is an admission you did not search. A gap is a specific
    capability that existing methods provably lack, or a specific condition
    under which they measurably fail.
\end{itemize}

\lead{Contributions: at most four, each refutable.} A contribution is
something a reviewer could disagree with. ``We propose a novel framework''
cannot be disagreed with, so it contributes nothing. Prefer artifacts and
claims: a model, a law, a dataset, a measured result, a proof. Give each a
forward reference to the section that delivers it --- this is a promise, and
the reader will check.

\lead{State what you do not claim.} One or two sentences saying explicitly
what is outside the paper's scope buys enormous credibility and pre-empts the
review that judges you against a paper you did not write.

% =========================================================================
\section{Related work}

Shorter than you think --- often three quarters of a column, sometimes folded
into the introduction and omitted entirely.

\begin{itemize}
  \item \textbf{One contrast per paragraph.} Each paragraph names a family of
    prior work and says what it does that you do differently, in one sentence.
    A list of summaries with no contrast is a literature review, not a related
    work section.
  \item \textbf{Do not explain trade-offs the reader has no context for.}
    Related work placed early can only be understood in terms of your method,
    which the reader has not read yet. Keep it comparative and brief; put the
    detailed positioning in the discussion where it will land.
  \item \textbf{Cite the work you actually build on more prominently than the
    work you merely list.}
\end{itemize}

% =========================================================================
\section{Method}

The largest section, and the one most often written as a sequence of
implementation details rather than an argument. It should read as one
mechanism, presented in stages.

\lead{Order: intuition, then formalism, then algorithm, then cost.} Give the
reader the idea in prose before the notation. A reader who understands the
idea will forgive dense notation; a reader who meets the notation first will
not reach the idea.

\lead{Every subsection is a stage of one mechanism.} If a subsection could be
removed and the method would still work, it belongs in an ablation, not the
method. If two subsections do unrelated things, the paper has two ideas and
should probably have one.

\lead{Teach only the background you use.} A background subsection that covers
the field is padding. Cover exactly the tools the method combines, in the
notation the method will use, and stop.

\subsection*{Notation}

\begin{itemize}
  \item Define every symbol once, at first use, and never redefine it.
  \item Fewer symbols is better. If a quantity appears twice, name it in
    words instead.
  \item Keep a consistent convention (vectors, matrices, sets, indices) and
    state it if it is not obvious.
  \item Equations are part of the sentence. Punctuate them, and make sure the
    text around them says what they mean --- an equation that is never
    referred to in prose is decoration.
  \item Number only the equations you refer to.
\end{itemize}

\subsection*{What the method section owes the reader}

\begin{itemize}
  \item \textbf{Reproducibility.} Enough that a competent reader could
    reimplement it. Where a value was chosen rather than derived, say so and
    say how it was chosen.
  \item \textbf{Cost.} The computational and memory cost of each operation,
    in the units that matter for the claim. If the paper's point is
    efficiency, this is not a detail --- it is the result.
  \item \textbf{Design justification, not just description.} For each
    non-obvious choice, one clause on why the alternative was rejected. This
    is where you pre-empt half the review.
  \item \textbf{Honest assumptions.} State the conditions under which the
    method is valid, in the method section, before the results depend on them.
\end{itemize}

\lead{Prefer a worked example to another paragraph of abstraction.} One
concrete instance --- a single observation entering the system, a single query
answered --- does more for comprehension than a page of generality.

% =========================================================================
\section{Results}
\label{sec:results}

The results section is not a report of what you ran. It is a sequence of
claims, each with evidence attached.

\lead{Declare the questions, then close them.} Open the section with the
specific questions the experiments answer, and let the subsection titles carry
them. Close each subsection by saying which question it settled and how. A
reader should be able to read only the first and last sentence of each
subsection and get the argument.

\lead{Claim-first paragraphs.} The first sentence of each results paragraph
states the finding; the rest supports it. Do not narrate the experiment and
reveal the conclusion at the end --- reviewers skim, and a paragraph whose
point arrives last reads as having no point.

\subsection*{How to state a number}

Every quantitative claim carries four things: the \textbf{number}, the
\textbf{named baseline}, the \textbf{named condition}, and a \textbf{pointer}
to the table or figure. ``Ours is better'' is not a claim. ``0.36\,m median
error against 0.29\,m for exact nearest neighbour on the same descriptors
(Table~II)'' is.

\begin{itemize}
  \item \textbf{Freeze numbers verbatim.} The same figure appears identically
    in the abstract, the introduction, the table and the conclusion. Rounding
    drift between sections reads as carelessness and invites the reviewer to
    check everything else.
  \item \textbf{Give errors a reference point.} An error means nothing
    without the floor it approaches or the cost it bought. Report it against
    the sensor's resolution, the annotation precision, human performance, or
    the resource it was traded for --- in the same sentence.
  \item \textbf{Denominate in the constrained resource.} If the paper is
    about a budget, price every result in that budget's units and against a
    hard limit. Put the resource in its own table column so it cannot be
    overlooked.
  \item \textbf{Report variance, not just means.} A single run is an
    anecdote. Give the spread, the number of seeds or trials, and enough for
    the reader to judge whether a difference is real.
  \item \textbf{Do not report more precision than you measured.} Three
    decimal places on a five-trial mean is a claim about noise you have not
    earned.
\end{itemize}

\subsection*{Baselines}

\begin{itemize}
  \item \textbf{Compare against the strongest version of the alternative.}
    Re-tune its hyperparameters, match its operating point, give it the
    privileged information it would have in its own paper.
  \item \textbf{Reproduce the published result before beating it.} If your
    measurement of a competitor differs from their reported number, resolve
    that first --- and then say you did. It is one of the most persuasive
    sentences available, because almost nobody does it.
  \item \textbf{Say what you strengthened, explicitly.} Steelmanning that the
    reader has to infer is steelmanning you get no credit for.
  \item \textbf{Match what varies.} If the comparison changes two things at
    once, it measures nothing. Hold the frontend, the data, and the metric
    fixed, and say that you did.
\end{itemize}

\subsection*{Ablations}

Each ablation isolates one variable and answers one question: \emph{is this
component doing what we claim?} Two useful disciplines:

\begin{itemize}
  \item \textbf{Include the control that could embarrass you.} If a component
    might be replaceable by something trivial --- a random choice, a constant,
    a simpler heuristic --- run that comparison. If the trivial version wins,
    you have learned something important before a reviewer does.
  \item \textbf{Run the ablation that makes the method look worst}, and
    explain the result. Papers that do this are believed about everything
    else.
\end{itemize}

\subsection*{Losses and negative results}

Put them in the main table, with a one-sentence mechanistic diagnosis in the
results text and the mechanism echoed in the discussion. A loss you explain
becomes evidence that you understand the system; a loss the reviewer finds
becomes evidence that you did not. If the loss is consistent with your thesis
--- you traded accuracy for size, and here is the accuracy --- then reporting
it \emph{is} the argument, not a concession to it.

\lead{Diagnose, do not merely disclose.} Leaving a lost row in a table
unremarked is survivable but wasteful: it reads as an oversight rather than a
finding. Naming the mechanism converts the same number into a demonstration
that you know where your method's boundary is, and pre-empts the reviewer who
was about to point it out.

\lead{Placement matters more than candour.} Be candid in the results and the
discussion, where a reader has the context to weigh the loss. Do not lead the
abstract with a deficit: state the trade axis and its scope condition there
instead. Honesty about a boundary reads as command of the problem; honesty
about a shortfall, before the reader knows what was bought with it, reads as
weakness.

Distinguish clearly between three different statements, because they carry
different weight: \emph{we measured no effect} (with the power to detect one),
\emph{we could not measure it} (underpowered or confounded), and \emph{we did
not test it}. Papers lose credibility by collapsing the second and third into
the first.

% =========================================================================
\section{Figures and tables}

Figures and tables are read before the text and often instead of it. Treat
each as a small paper.

\begin{itemize}
  \item \textbf{One message per figure.} If you cannot say what a figure
    proves in a sentence, it is two figures or none.
  \item \textbf{Captions are self-contained and state the takeaway.} Forty to
    eighty-five words, present tense, opening with the claim rather than a
    description of the panels. A reader who reads only your figures and
    captions should get the paper's argument.
  \item \textbf{Put the numbers in the caption.} Normalised into units a
    reader feels --- a distance as a fraction of the robot's height, a
    footprint as a ratio to the baseline it replaces.
  \item \textbf{The first figure shows a result, not an architecture.} The
    strongest opening figures show failure-then-fix, or the system's own
    output. A block diagram tells the reader how you built it, which is not
    yet a reason to care.
  \item \textbf{Answer one objection in the caption.} A single clause
    pre-empting the obvious criticism is among the cheapest credibility moves
    available.
  \item \textbf{Axes labelled, with units.} Legends that do not require
    cross-referencing the text. Consistent colours across figures, and check
    that they survive greyscale printing.
\end{itemize}

\lead{Table design.} What you vary goes in rows; what you measure goes in
columns. Bold the best result per column and say in the caption what bold
means. Include the resource column (memory, time, parameters) even when it is
not the headline --- it is what a sceptical reader looks for. Align decimals.
Do not make the reader compute a ratio you could have given them.

% =========================================================================
\section{Limitations and conclusion}

Merge limitations into the discussion rather than giving them their own
heading, and give each one a number and a path forward. ``Our method may not
generalise'' is worthless; ``performance degrades beyond room scale, where the
map exceeds $N$ items; extending to building scale requires $X$'' is a
research direction and reads as command of the problem.

The conclusion restates the claim, states the trade plainly, and stops. It is
not a summary of the paper --- the reader has just read the paper.

% =========================================================================
\section{Sentences and paragraphs}

\begin{itemize}
  \item \textbf{Topic sentence carries the claim.} One idea per paragraph,
    announced first.
  \item \textbf{Active voice for what you did.} ``We measured'' rather than
    ``it was measured''. Passive is for when the actor genuinely does not
    matter.
  \item \textbf{Verbs must carry mechanism.} \emph{Provides}, \emph{enables},
    \emph{allows}, \emph{leverages} and \emph{utilises} say nothing. Say what
    the thing does.
  \item \textbf{Cut throat-clearing.} ``It is worth noting that'', ``In order
    to'', ``It should be mentioned'', ``As can be seen from''. Delete and the
    sentence improves.
  \item \textbf{No promotional adjectives.} \emph{Novel}, \emph{powerful},
    \emph{elegant}, \emph{extensive}. Let the result be impressive; do not
    assert that it is.
  \item \textbf{Hedge mechanism and generality; never hedge measurements.}
    State what you measured flatly. Mark speculation explicitly as
    speculation. A paper that hedges its data reads as unsure the data is
    right.
  \item \textbf{Prefer the specific word.} ``Improves accuracy'' is weaker
    than ``raises recall at one metre from 0.31 to 0.44''.
  \item \textbf{Same thing, same name.} Do not elegantly vary your
    terminology; a reader cannot tell whether a new word means a new concept.
\end{itemize}

% =========================================================================
\section{Citations, code and the red-team pass}

\subsection*{Citation hygiene}

\begin{itemize}
  \item \textbf{Cite the source, not the survey.} A survey is the right
    citation for a landscape claim and the wrong one for a specific result.
    Citing a review for something a particular paper proved is the tell of
    someone who read neither.
  \item \textbf{Never cite a paper for something it does not say.} Reviewers
    are frequently the authors of the work you are citing, and misattribution
    is one of the few errors that reliably converts a positive review into a
    hostile one. If you are characterising someone's result, open their paper
    and check the sentence.
  \item \textbf{Cite the work you build on prominently, and the work you
    merely acknowledge briefly.} A flat wall of thirty citations tells the
    reader nothing about what mattered.
  \item \textbf{Self-cite in the third person under anonymous review.} Write
    ``the approach of [12]'' rather than ``our previous work [12]'', and check
    that the surrounding sentence does not deanonymise you anyway. Do not
    remove the citations --- omitting relevant prior work is worse than the
    small risk of being guessed.
  \item \textbf{Fix the bibliography.} Wrong venues, missing years, arXiv
    entries for papers that were published properly, inconsistent
    capitalisation --- individually trivial, collectively an impression of
    carelessness formed before the reviewer reaches your results.
\end{itemize}

\subsection*{Code and reproducibility}

Reviewers rarely run code, but its presence and shape change how claims are
read. An anonymised repository, or a clear statement that code will be
released, converts ``we implemented a fair version of the baseline'' from an
assertion into a checkable one.

\begin{itemize}
  \item Under anonymous review, use an anonymised repository service; a
    personal or lab account link is a deanonymisation.
  \item Give the reader what makes the numbers reproducible: seeds, the
    hyperparameters you actually used, dataset versions and splits, and the
    hardware the timings were measured on. A hyperparameter table costs a
    quarter column and removes an entire class of objection.
  \item If you cannot release code, say why, and compensate with detail ---
    an algorithm box and complete settings.
\end{itemize}

\subsection*{The red-team pass}

A week before the deadline, give the draft to someone --- ideally not a
co-author --- with one instruction: \emph{find the reason to reject this}.
Not ``what do you think'', which produces politeness, but a mandate to
attack. Ask them specifically for: the claim they cannot verify from the
paper alone; the baseline they would have wanted; the sentence they did not
believe; and the first paragraph where they lost the thread.

Then triage. Objections that can be answered by a sentence, a column, or a
caption should be answered before submission --- that is the whole point of
doing it a week early. Objections that need an experiment you cannot run
become an explicit limitation, which is far better than leaving a reviewer to
discover them. If you have no willing reader, the substitute is to leave the
draft for two days and read it once, quickly, in the role of someone who
wants to reject it; this is much weaker than a real reader, but it is not
nothing.

% =========================================================================
\section{Getting accepted in a track}

Everything above helps a paper be admired. Acceptance is a different and
earlier problem, decided by a handful of anonymous reviewers working from a
scoring form, and the writing that survives that process is not quite the
writing that wins prizes. Prize-winning prose is bold and coins a frame;
accepted prose is bold \emph{and} defensive --- it closes the objection
before the reviewer has to raise it. A student's first paper needs the second
skill more than the first.

\lead{Your first keyword chooses your reviewers.} At IEEE conferences the
author's first keyword routes the paper to an Editor, who selects the
Associate Editor, who recruits the reviewers --- and scores are then
normalised against the other papers filed under that keyword. That single
dropdown therefore decides both who judges the work and what it is judged
against. Choose the community that will recognise your contribution
\emph{as} a contribution: a paper about memory efficiency filed under a
keyword whose other submissions are all accuracy papers will be read by
people for whom your headline result is a loss. This is the highest-leverage
thirty seconds of the whole submission.

\lead{Assume there is no second turn.} Some venues run an author-response
phase; others do not, and at those the submitted PDF is the only thing that
will ever argue your case. Check the call for papers, and if there is no
rebuttal, write as though every reviewer objection must be answered inside
the paper. This is what makes the in-text objection-answering habit
(Section~\ref{sec:results}) a scoring matter rather than a stylistic one.

\subsection*{Learn the conventions from the tables, not the prose}

The single most useful thing you can do before designing an experiment:
\textbf{take the five papers whose rows you intend your results to sit
beside, and read their tables --- the headers, the units, the scene lists,
the columns that always exist --- then adopt those axes exactly.}

A track's real expectations live in its table headers. If the field reports
median error in centimetres over all seven scenes of a benchmark, and you
report mean error in metres on one scene, your number is not worse --- it is
\emph{uncomparable}, and an uncomparable number cannot be credited however
good it is. Reviewers experience this as an inability to place you, which
converts into a middling score rather than an argument.

Two corollaries:

\begin{itemize}
  \item \textbf{Claiming a benchmark's axis obliges you to cite the line of
    work that owns it.} Reporting on a standard dataset while ignoring the
    method family that defined its numbers reads as unfamiliarity with the
    field, which is more damaging than a weak result.
  \item \textbf{Report the whole benchmark, not the convenient subset.} One
    scene, one metric, or one parameter setting is not a result, and a
    partial table invites the reviewer to assume the omitted cases are worse.
\end{itemize}

\subsection*{Defuse your subfield's standing objections}

Every subfield has three or four objections it raises reflexively. They are
discoverable: read recent accepted papers and notice what they pre-empt
without being asked. The presence of a defusing move is visible even when the
objection is not. Then perform the same move, explicitly. The general
patterns:

\begin{itemize}
  \item \textbf{``Compared against what, in what units?''} Any efficiency or
    compactness claim belongs as a \emph{column in the same table} as the
    accuracy number, on the same data, beside the baselines --- never as an
    adjective in prose. And if you claim a resource, you owe its companion:
    a memory claim without a latency number invites the suspicion that the
    latency is bad.
  \item \textbf{``It only works at the scale you tested.''} Give one
    stress axis --- performance against scene size, object count, sequence
    length --- even if it degrades. A measured degradation curve reads as
    characterisation; silence reads as a limit you are hiding.
  \item \textbf{``The gain comes from a component, not your contribution.''}
    Hold everything else fixed and vary only your part. If your system
    consumes another system's outputs, run both backends on the identical
    stream and say so in the abstract, not in a mid-section aside.
  \item \textbf{``No real-robot results.''} If your named comparison has
    hardware and you do not, the asymmetry is visible. Short of a robot, use
    real recorded sensor data rather than simulation alone, and use the video
    attachment.
\end{itemize}

\subsection*{Do not import another venue's shape}

Venues differ in ways that are invisible until they cost you. Page limits may
or may not include references; supplementary material may be unlimited,
capped, or forbidden; a Limitations section is mandatory at some venues and
merely optional at others; review may be single- or double-anonymous. A paper
shaped for a venue that excludes references from its limit and permits a large
appendix can carry one to two pages more content than the same work at a venue
that counts everything --- so a structure copied across without checking will
either overflow or waste space.

Where a video is permitted it is often the only overflow channel available.
Make it self-contained and free of identifying branding, and make sure the
paper stands without it: reviewers are not always required to watch.

\subsection*{What a score means}

Review scales run from a clear accept, through a wide borderline band, to a
clear reject. Most papers land in the borderline band, and most decisions are
made there --- which means the marginal effort is best spent not on making a
good result better, but on removing the specific reasons a sympathetic
reviewer would still hesitate: an uncomparable metric, a missing baseline, an
unexplained loss, a claim without a condition, a limitation the reviewer
found before you admitted it.

% =========================================================================
\section{After submission}

\lead{Reading reviews.} Reviews arrive blunt, sometimes wrong, occasionally
both. Read them once, then leave them for a day before deciding anything ---
the first reading is about your feelings and the second is about the paper.
Then sort each point into: correct and fixable; correct and not fixable now;
a misunderstanding; and plainly wrong.

The most useful reframing is that a misunderstanding is usually a writing
defect. If a competent reviewer misread your claim, other readers would have
too, and the sentence that misled them is findable. Reviewers who
misunderstand are giving you free evidence about where the paper is unclear,
which is worth more than the score.

\lead{Rebuttals, where they exist.} Some venues invite an author response;
where they do, it is short and read quickly, so it must be triaged rather
than complete.

\begin{itemize}
  \item Lead with the factual corrections that change a score --- a
    misreading of what you claimed, a baseline the reviewer thought was
    missing that is in Table~III --- with a pointer to where the answer
    already is.
  \item Concede the fair points plainly and say what you will change. A short
    concession buys credibility for the paragraph where you disagree.
  \item Answer what can be answered with evidence you already have. Promising
    experiments you have not run is transparent and weakens everything else.
  \item Do not argue tone, do not relitigate taste, and do not spend space on
    the reviewer you will never persuade. Write for the borderline reviewer
    and the AE reading the thread.
\end{itemize}

\lead{Rejection.} Most papers are rejected at least once, including good
ones, and the reviews you get are an asset the next submission does not have
to earn. Turn them into a checklist, fix what is fixable, and resubmit
promptly --- the same work is often a strong submission to a sibling venue or
journal within a couple of months. Two habits make this cheap: keep the
reviews and your responses in the repository beside the paper, and fix the
writing defects immediately while you can still remember what you meant.

The failure mode to avoid is treating a rejection as a verdict on the
research rather than on this presentation of it. It is nearly always the
second.

% =========================================================================
\section{Anti-patterns}

\begin{itemize}
  \item A title with three claims, or an opener from the banned list.
  \item Five or more metrics scattered through the abstract.
  \item An abstract ending on the effectiveness of the proposed method.
  \item ``Little work has been done'' as a gap statement.
  \item Contributions that cannot be disagreed with.
  \item A background section covering the field rather than the method.
  \item Terse captions (``System overview.'').
  \item A bare architecture diagram as the first figure.
  \item Results paragraphs that narrate the experiment and conclude at the
    end.
  \item Bare percentages with no baseline, condition, or pointer.
  \item Multiple vague ``first'' claims. Make at most one, tightly scoped.
  \item Hiding a loss where a reviewer will find it.
  \item Anything load-bearing in supplementary material.
\end{itemize}

% =========================================================================
\section{Revision checklist}

Run against a draft that already exists, roughly in order of value per minute.

\begin{itemize}[leftmargin=*]
  \ck The claim is one falsifiable sentence, and it appears in the same words
    in title, abstract, body and conclusion.
  \ck Abstract is seven sentences; every sentence loses information if
    deleted; one headline number; the central result is present; it does not
    end on effectiveness.
  \ck Title: one claim, under about eleven words, no banned openers.
  \ck Contributions: at most four, each refutable, each with a forward
    reference that the paper honours.
  \ck Every headline number identical everywhere it appears.
  \ck Every quantitative claim has number, baseline, condition, pointer.
  \ck Errors reported against a floor or an exchange rate.
  \ck Variance and trial counts reported.
  \ck Baselines steelmanned, and the steelmanning stated explicitly.
  \ck Losses in the main tables, each with a mechanistic diagnosis.
  \ck Method reads as one mechanism in stages; background covers only what is
    used; costs stated.
  \ck Every symbol defined once; every numbered equation referenced in prose.
  \ck First figure shows a result; caption states the claim, carries numbers,
    pre-empts one objection.
  \ck Every figure has one message; every caption is self-contained.
  \ck Tables have a resource column and a stated bolding convention.
  \ck Limitations merged into discussion, quantified, with a path.
  \ck No promotional adjectives; no mechanism-free verbs; no throat-clearing.
\end{itemize}

\noindent Before submitting:

\begin{itemize}[leftmargin=*]
  \ck Someone was asked to reject the draft, and their answerable objections
    were answered.
  \ck Citations checked: sources not surveys, nothing cited for a claim it
    does not make, self-citations in the third person, bibliography clean.
  \ck Seeds, hyperparameters, splits and timing hardware stated; code link
    anonymised if review is anonymous.
  \ck Metrics, units and dataset splits match the papers your results sit
    beside --- checked against their tables, not their prose.
  \ck The benchmark is reported whole, not as a convenient subset.
  \ck The work that owns each benchmark you use is cited and, where
    relevant, compared against.
  \ck Your subfield's three standing objections are each answered in the
    text, unprompted.
  \ck Any resource claim is a table column, with its companion measurement
    (memory with latency, speed with accuracy).
  \ck One stress axis reported, even where performance degrades.
  \ck First keyword chosen for the community that will recognise the
    contribution.
  \ck Every reviewer objection you can predict is answered inside the paper,
    on the assumption that there is no rebuttal.
\end{itemize}

% =========================================================================
\section*{Before submission}
\addcontentsline{toc}{section}{Before submission}

Check the current call for papers rather than trusting a previous cycle, or
a colleague's memory of one:

\begin{itemize}
  \item the page limit, and whether references count toward it;
  \item whether review is anonymous, and what must therefore be stripped from
    the author block, acknowledgements, grant numbers, figure captions and
    self-citations;
  \item whether supplementary material is permitted and whether reviewers are
    obliged to look at it, or whether everything load-bearing must sit inside
    the page limit;
  \item whether there is an author-response phase, since its absence means
    every objection must be pre-empted in the submitted PDF;
  \item video specifications, if a video is allowed;
  \item any disclosure requirements, including for AI tool use;
  \item whether a preprint is permitted, and whether a prior version carrying
    a DOI makes the work ineligible.
\end{itemize}

\end{document}

